% ==============================================================================
% preambulo.tex
% Preámbulo común de las dos presentaciones (best_equipo_01 y compare_equipo_01).
% Las figuras y tablas se toman de 03_Output, donde las generan los scripts del
% repositorio, y las cifras citadas vienen de 03_Output/Tables/pres_cifras.tex
% (01_Scripts/04_Results/03_figuras_presentaciones.R).
% ==============================================================================
\usepackage[T1]{fontenc}
\usepackage[utf8]{inputenc}
\usepackage{lmodern}
\usepackage{graphicx}
\usepackage{booktabs}
\usepackage{tikz}
\usetikzlibrary{positioning, arrows.meta}

\graphicspath{{../03_Output/Figures/}}
\newcommand{\maeprivado}{151,5}
\newcommand{\preciotipico}{806}
\newcommand{\maepctprecio}{19\%}
\newcommand{\ntrain}{36.407}
\newcommand{\nrawtrain}{38.644}
\newcommand{\ncopias}{2.237}
\newcommand{\ntest}{10.286}
\newcommand{\nchaptrain}{305}
\newcommand{\pcttestchap}{94\%}
\newcommand{\preciomdchap}{750}
\newcommand{\preciomdresto}{550}
\newcommand{\aptotest}{97\%}
\newcommand{\aptotrain}{75\%}
\newcommand{\sinareatest}{44\%}
\newcommand{\sinareachap}{6\%}
\newcommand{\coordredondeada}{52\%}
\newcommand{\coordredtrain}{40\%}
\newcommand{\nvarexternas}{61}
\newcommand{\nvartexto}{44}
\newcommand{\ntopicos}{15}
\newcommand{\nenvios}{25}
\newcommand{\margenchap}{19\%}
\newcommand{\sobreestimachap}{46\%}
\newcommand{\sobrediezchap}{25\%}
\newcommand{\pctverde}{11\%}
\newcommand{\pctamarillo}{61\%}
\newcommand{\pctrojo}{27\%}
\newcommand{\mejormargen}{17\%}
\newcommand{\gananciacien}{1.153}
\newcommand{\negociosmargen}{42\%}
\newcommand{\descuentomax}{15\%}
\newcommand{\pasoentorno}{96}
\newcommand{\pasoxgb}{37}
\newcommand{\pasotexto}{13}
\newcommand{\ventajaxgb}{39}
\newcommand{\ventajaensamble}{3}
\newcommand{\mismabasedif}{58}
\newcommand{\cartdistintos}{59}
\newcommand{\sesgobarato}{17\%}
\newcommand{\sesgocaro}{19\%}


% Colores
\definecolor{azul}{HTML}{2166AC}
\definecolor{rojo}{HTML}{B2182B}
\definecolor{verde}{HTML}{1A9850}
\definecolor{gris}{HTML}{595959}
\definecolor{grisclaro}{HTML}{F2F2F2}

% Tema sobrio: títulos en azul, sin íconos de navegación, número de página
\usetheme{default}
\setbeamertemplate{navigation symbols}{}
\setbeamercolor{frametitle}{fg=azul}
\setbeamerfont{frametitle}{series=\bfseries, size=\large}
\setbeamercolor{title}{fg=azul}
\setbeamerfont{title}{series=\bfseries}
\setbeamercolor{structure}{fg=azul}
\setbeamertemplate{itemize items}[circle]
\setbeamertemplate{footline}{%
  \hfill\textcolor{gris}{\scriptsize Equipo 1 \quad \insertframenumber/\inserttotalframenumber}\hspace{1em}\vspace{0.6em}}

% Cifra grande con su descripción. El argumento opcional es el ancho (por
% defecto, un tercio del ancho de la diapositiva).
\newcommand{\cifra}[3][0.3\textwidth]{%
  \begin{minipage}[t]{#1}\centering
    {\Huge\bfseries\color{azul}#2}\\[0.3em]
    {\small\color{gris}#3}
  \end{minipage}}

% Cifra mediana, para acompañar un mapa sin quitarle espacio
\newcommand{\cifrachica}[2]{%
  {\LARGE\bfseries\color{azul}#1}\\[0.1em]
  {\footnotesize\color{gris}#2}}

% Mensaje corto debajo de una figura
\newcommand{\lectura}[1]{\par\vspace{0.2em}{\scriptsize\color{gris}#1}}

% Diapositivas de mapa vertical: el mapa ocupa toda la altura a la izquierda y
% el título va arriba de la columna derecha, con el mismo estilo de los demás
% títulos. \altomapa es la altura del mapa.
\newcommand{\altomapa}{1.02\textheight}
\newcommand{\titulolateral}[1]{%
  {\usebeamerfont{frametitle}\usebeamercolor[fg]{frametitle}#1\par}\vspace{0.9em}}
\newcommand{\mapavertical}[1]{\vspace*{-0.4em}\includegraphics[height=\altomapa]{#1}}

% Figura ancha que usa casi todo el ancho de la diapositiva
\newcommand{\figuraancha}[2][0.97]{\makebox[\textwidth]{\includegraphics[width=#1\paperwidth]{#2}}}
